\documentclass[10pt,a4paper]{amsart}
\usepackage[margin=19mm]{geometry}
\usepackage{amsmath,amssymb,mathtools}
\usepackage[scr=rsfs,cal=euler]{mathalfa}
\usepackage{xfrac}
\usepackage[unicode,hidelinks,bookmarks=false]{hyperref}
\newcommand{\ie}{i.e.\ }
\newcommand{\cf}{cf.\ }
\newcommand{\resp}{respectively\ }
\newcommand{\Subsection}{\subsection}
\providecommand{\nobreakdash}{\nobreak\hbox{-}}
\setlength{\emergencystretch}{2em}
\multlinegap0pt
\usepackage{bbm}
\usepackage{amssymb}
\usepackage{bm}
\usepackage{cancel}
\usepackage{enumerate}
\usepackage{tensor}
\newtheorem{proposition}{Proposition}
\newcommand{\R}{\mathbb{R}}
\newcommand{\overbar}[1]{\mkern 1.5mu\overline{\mkern-1.5mu#1\mkern-1.5mu}\mkern 1.5mu}
\title{\Huge{Nonlinear diffusion in\\ relativistic kinetic theory}}
\author{Simone Calogero}
\date{Source: arXiv:2508.20147v2 (CC BY 4.0)}
\begin{document}
\maketitle

\noindent\emph{Source and licence.} \Huge{Nonlinear diffusion in\\ relativistic kinetic theory}. Version arXiv:2508.20147v2 (CC BY 4.0), distributed by arXiv under the Creative Commons Attribution 4.0 International licence, http://creativecommons.org/licenses/by/4.0/, as recorded in the arXiv licence metadata for this exact version and shown on the abstract page https://arxiv.org/abs/2508.20147v2. Full licence notice: \texttt{licenses/arxiv-2508.20147-CC-BY-4.0.txt}.\par\smallskip

\noindent\textit{Editorial scope.} Counted unit: the article's proposition inv4dprop (source0.tex lines 373--375). The article's complete original proof of it is delivered with it (source0.tex lines 376--396, one \texttt{proof} environment of 21 lines). No mathematics is written, repaired or re-derived: the statement and the proof are the article's own lines, in the article's own notation, and no retained line is substituted or re-flowed.  Retained uncounted as the source-local context the proof needs: lines 257--259; lines 269--272; lines 275--277; lines 301--303; lines 312--314.  Nothing was omitted from any retained span, so the package carries no omission record.  No retained line contains a citation, so the wrapper carries no bibliography and no bibliography entry.  No retained line contains a figure environment, an image-inclusion command or drawing code, so no image is delivered and the package declares no asset dependency.  Editorial formatting only, in addition: an AMS \texttt{amsart} preamble that restates the article's own declarations and macros used by the retained lines, namely the environments \texttt{proposition} on separate counters, together with the standard packages those lines need.  The retained environments print in this document as Proposition 1 in order of appearance, whereas the article prints the same items with the numbering of its own full document.  The complete native source of the article as distributed in the arXiv e-print for this exact version is bundled unchanged as \texttt{sources/184005/source0.tex}.
\begin{equation}\label{eqf*}
u^\alpha\partial_{x^\alpha}\mathbbm{f}=\Phi_\mathbbm{f} \, \eta^{\alpha\beta}\partial_{u^\alpha}\partial_{u^\beta} \mathbbm{f}+\partial_{u^\alpha} (K_\mathbbm{f}^\alpha \, \mathbbm{f}),
\end{equation}

\begin{equation}\label{lambdaid}
\Lambda({\bm\xi})^T\eta\, \Lambda({\bm \xi})=\eta,
\quad \Lambda({\bm \xi})^{-1}=\Lambda(-{\bm \xi}),
\end{equation}

\begin{equation}\label{transphik}
\Phi_{\bar{\mathbbm{f}}}(x)=\Phi_\mathbbm{f} (\bar{x}),\quad {K}_{\bar{\mathbbm{f}}}(x,u) =\Lambda(-{\bm \xi}) K_\mathbbm{f} (\bar{x},\bar{u}).
\end{equation}

\begin{equation}\label{eqf2}
u^\alpha\partial_{x^\alpha}f=u^0\partial_{u^i}\left(\Phi\frac{(h^{-1})^{ij}}{u^0}\partial_{u^j}f+\frac{K^i}{u^0}f\right),
\end{equation}

\begin{equation}\label{N}
N^\alpha(x) =\int u^\alpha f(x,{\bm u})\frac{d{\bm u}}{u^0}
\end{equation}

\begin{proposition}\label{inv4dprop}
$\Phi$, $K$ satisfy the transformation laws~\eqref{transphik}; in particular,~\eqref{eqf*}, and thus also~\eqref{eqf2}, are invariant under Lorentz transformations.
\end{proposition}

\begin{proof}
To prove the claim, we rewrite $\Phi$ and $K$ in terms of $\mathbbm{f}$ (instead of $f$). The definition~\eqref{N} of $N$ is equivalent to 
\[
N(x)=\int_\mathcal{U} u\,\mathbbm{f}(x, u)\,du,
\]
where $\mathcal{U}=\{u\in \R^4:u^0=\sqrt{c^2+|{\bm u}|^2}\}$. By the change of variable $u=\bar{v}=\Lambda({\bm\xi}) v$ we obtain
\[
N(\bar{x})=\int_\mathcal{U} \Lambda({\bm\xi}) v\,\mathbbm{f}(\bar{x}, \bar{v})\,dv=\Lambda({\bm\xi})\int_\mathcal{U} v\,\bar{\mathbbm{f}}(x,v)\,dv:=\Lambda({\bm\xi}) \overbar{N}(x),
\] 
where we used that $\det\Lambda({\bm\xi})=1$, and that $\mathcal{U}$ is invariant under Lorentz transformations.
Hence, using~\eqref{lambdaid},
\begin{align*}
K_\mathbbm{f} (\bar{x},\bar{u})&=\mu \left(\bar{u}+\frac{c^2N(\bar{x})}{N(\bar{x})^\alpha \bar{u}_\alpha}\right)=\mu\Lambda({\bm\xi})\left(u+\frac{c^2\overbar{N}(x)}{\overbar{N}(x)^\alpha u_\alpha}\right)=\Lambda({\bm\xi})K_{\bar{\mathbbm{f}}}(x,u)\\
&\Rightarrow K_{\bar{\mathbbm{f}}}(x,u)=\Lambda(-{\bm \xi})K_\mathbbm{f} (\bar{x},\bar{u}).
\end{align*}
The claim on $\Phi$ is proved likewise:
\begin{align*}
\Phi_\mathbbm{f} (\bar{x})&=\frac{c^2\mu}{3}\left(1+\int_\mathcal{U} \frac{c^2\mathbbm{f}(\bar{x},u)}{N(\bar{x})^\alpha u_\alpha}\,du\right)=\frac{c^2\mu}{3}\left(1+\int_\mathcal{U} \frac{c^2\mathbbm{f}(\bar{x},\bar{v})}{N(\bar{x})^\alpha \bar{v}_\alpha}\,d v\right)\\
&=\frac{c^2\mu}{3}\left(1+\int_\mathcal{U} \frac{c^2\bar{\mathbbm{f}}(x,v)}{\overbar{N}(x)^\alpha v_\alpha}\,dv\right)=\Phi_{\bar{\mathbbm{f}}}(x).
\end{align*}
\end{proof}

\end{document}
